% !TEX root = ../main.tex
\section*{Discussion}

The weighted statistic changes one step of SatuTe. The branch, subtree likelihood vectors, null hypothesis and one-sided test remain the same. When the fitted substitution model separates the other non-stationary components, the statistic includes them after scaling by their expected decay along branch \(AB\).

The simulations identify the settings in which the formula changes the decision. Under JC, the statistic is unchanged. Under four-category rate-heterogeneous nucleotide and protein models, the weighted statistic changed decisions most often for 100-site alignments and long tested branches. With more sites, the published SatuTe statistic usually had enough information to call the branch phylogenetically informative, and the weighted statistic rarely changed the binary decision. Across all 1,728 matched-protein cells the observed change was never negative, but its size varied markedly among LG, WAG, JTT and Q.pfam. The independent truth-labelled experiment found nearly identical empirical FDR and power for the two formulas away from this transition. The evidence therefore supports a localized gain at the saturation decision boundary, not a universal increase in power under every model and branch regime.

The explicit invariant-site analysis identifies a separate issue. A zero-rate category never loses the shared ancestral state across the focal branch, so treating that category as independent in the infinite-branch null creates a positive score even when the data are generated from the correct null. Studentization then amplifies the shift at rate \(\sqrt n\), explaining why the counterfactual nominal rejection rate approaches one as alignment length grows. The refined denominator restores exact null centering and ordinary standard-normal calibration. Once both formulas are given empirical thresholds, the remaining power change is smaller and localized to the detection boundary, where the refined statistic gained as much as 16.3 percentage points in the tested grid.

The dense grid also clarifies resolution. A coarse set of branch lengths captures the false-positive failure but understates the peak power difference because the detection transition moves with alignment length and rate mixture. The 2.5-unit steps from 10 to 25 and 5-unit steps from 30 to 50 resolve the diagonal power-gain ridge in the full-grid figure. Outside this ridge, both formulas usually have either near-null or near-complete power, so a formula change has little room to alter the binary decision.

The branch-neighborhood analysis addresses branch specificity. If the weighted statistic simply made all branches easier to call informative, the same change would appear on neighboring branches. In the 1,000-replicate simulations, the decision changes were concentrated on the tested branch, while nearest neighbors were already informative under both formulas. This is consistent with a branch-specific interpretation within the simulated five-taxon and 16-taxon settings.

In the Tree-of-Life \(+G4\) examples, the recomputations changed sliding-window saturation calls. These examples do not by themselves quantify the formula effect, because the published and native curves also differ in input files and implementation. The paired full-alignment simulations therefore provide the more controlled test. The sliding-window results are biological examples of the same question: whether regions of an alignment retain enough phylogenetic signal for a selected branch.

The original SatuTe paper did not resolve the true phylogenetic relationship of Archaea, Bacteria and Eukaryota, and this extension does not change that scope. SatuTe tests for saturation; it does not determine which tree is correct. The weighted statistic keeps that role. It applies when a rate-heterogeneous model has separated non-stationary eigenvalues and the analysis asks whether signal outside the second-largest-eigenvalue component should contribute. The published SatuTe coefficient remains the simpler reference statistic.
